% ─────────────────────────────────────────────────────────────────────────────
This appendix provides a full description of the BookForMe bilingual dataset,
including the annotation schema, sample utterances, class distribution, and
augmentation statistics.

% ─────────────────────────────────────────────────────────────────────────────
\section{Dataset Overview}
\label{app:data-overview}

The BookForMe dataset is a custom-built collection of 657 annotated utterances
representing real-world booking conversations for informal services in Karachi,
Pakistan.  It is the first publicly documented bilingual (English / Roman Urdu)
dataset for task-oriented booking dialogue in an informal South Asian service
context.

\begin{table}[H]
\centering
\caption{BookForMe dataset summary statistics.}
\label{tab:data-summary}
\begin{tabularx}{0.75\textwidth}{Xr}
\toprule
\textbf{Property} & \textbf{Value} \\
\midrule
Total utterances                & 657 \\
Training split (80\%)           & 526 \\
Test split (20\%)               & 131 \\
Number of intent classes        & 5 \\
Languages                       & English, Roman Urdu, Mixed (code-switched) \\
Service categories covered      & Futsal, Padel, Cricket, Salon, Gaming Zone \\
Collection period               & October 2025 -- January 2026 \\
Annotation tool                 & Google Sheets (manual) \\
Inter-annotator agreement       & Cohen's $\kappa = 0.87$ (substantial) \\
\bottomrule
\end{tabularx}
\end{table}

% ─────────────────────────────────────────────────────────────────────────────
\section{Annotation Schema}
\label{app:data-schema}

Each utterance in the dataset is annotated with three fields:

\subsection{Intent Labels}

\begin{table}[H]
\centering
\caption{Intent label definitions and examples.}
\label{tab:intent-defs}
\small
\begin{tabularx}{\textwidth}{@{}>{\raggedright\arraybackslash}p{0.20\textwidth}Y Y@{}}
\toprule
\textbf{Label} & \textbf{Description} & \textbf{Example Utterances} \\
\midrule
inquiry
  & User is checking availability or requesting to book a specific slot.
  & ``Is the 7~PM padel court available on Friday?'' /
    ``Mujhe kal shaam 6 baje futsal chahiye.'' \\[4pt]
info
  & User is asking for general information (price, amenities, location, hours).
  & ``What is the price per hour for the cricket net?'' /
    ``Parking hai yahan?'' \\[4pt]
transaction\_confirm
  & User is referencing a payment, confirming a transaction, or submitting proof.
  & ``I've sent the Rs 2000 advance.'' /
    ``Screenshot bhej diya hai.'' \\[4pt]
greeting
  & User is initiating conversation with a greeting or pleasantry.
  & ``Hi'' / ``Hello, are you there?'' / ``Aoa'' (Urdu greeting abbreviation) \\[4pt]
unknown
  & Out-of-scope, unintelligible, or irrelevant message.
  & ``I want pineapple juice.'' / ``???'' / Random keyboard mash \\
\bottomrule
\end{tabularx}
\end{table}

\subsection{Language Labels}

Each utterance is also tagged with one of three language labels:
\begin{itemize}
  \item \texttt{EN}: Purely English utterance.
  \item \texttt{UR}: Purely Roman Urdu utterance.
  \item \texttt{MIXED}: Code-switched utterance containing both English and Roman Urdu
        elements (the most common type in real vendor conversations).
\end{itemize}

\subsection{Slot Annotations}

For utterances where slots are extractable, a JSON dictionary of entities is
provided:

\begin{lstlisting}[style=jsonstyle, caption={Sample slot annotation for an inquiry utterance.}]
{
  "intent": "inquiry",
  "language": "MIXED",
  "text": "Mujhe Friday ko 7 baje ke baad paddle court chahiye under 6000",
  "slots": {
    "service": "padel court",
    "day_of_week": "Friday",
    "time_after": "19:00",
    "budget_max": 6000,
    "currency": "PKR"
  }
}
\end{lstlisting}

\begin{lstlisting}[style=jsonstyle, caption={Sample slot annotation for a transaction\_confirm utterance.}]
{
  "intent": "transaction_confirm",
  "language": "MIXED",
  "text": "Screenshot bhej diya hai. 2000 rupay transfer kiye.",
  "slots": {
    "payment_amount": 2000,
    "currency": "PKR",
    "payment_method": "transfer"
  }
}
\end{lstlisting}

% ─────────────────────────────────────────────────────────────────────────────
\section{Sample Utterances}
\label{app:data-samples}

Table~\ref{tab:data-samples} shows representative utterances from each intent class
in all three language varieties.

\begin{table}[H]
\centering
\caption{Sample utterances from the BookForMe dataset.}
\label{tab:data-samples}
\small
\begin{tabularx}{\textwidth}{>{\raggedright\arraybackslash}p{0.30\textwidth}p{0.10\textwidth}Y}
\toprule
\textbf{Intent} & \textbf{Lang.} & \textbf{Utterance} \\
\midrule
\texttt{inquiry} & EN
  & ``Do you have a futsal court available tomorrow at 9~PM?'' \\
\texttt{inquiry} & UR
  & ``Kal raat 9 baje futsal available hai?'' \\
\texttt{inquiry} & MIXED
  & ``Friday evening mein padel slot available hai under 1500?'' \\
\hline
\texttt{info} & EN
  & ``What's the price for a one-hour padel session?'' \\
\texttt{info} & UR
  & ``Kitna rate hai ek ghante ka?'' \\
\texttt{info} & MIXED
  & ``AC courts hain aur parking bhi hai kya?'' \\
\hline
\texttt{transaction\_confirm} & EN
  & ``I've sent the advance payment of Rs~2000 to the number.'' \\
\texttt{transaction\_confirm} & UR
  & ``Paisa bhej diya. Screenshot deta hoon.'' \\
\texttt{transaction\_confirm} & MIXED
  & ``2000 rupees transfer ho gayi. Screenshot attach kar raha hoon.'' \\
\hline
\texttt{greeting} & EN
  & ``Hello! Is anyone there?'' \\
\texttt{greeting} & UR
  & ``Aoa, koi hai?'' \\
\texttt{greeting} & MIXED
  & ``Hi, booking karni thi.'' \\
\hline
\texttt{unknown} & EN
  & ``Can you order me a pizza?'' \\
\texttt{unknown} & UR
  & ``Yaar mujhe neend aa rahi hai.'' \\
\texttt{unknown} & MIXED
  & ``asdfghjkl'' \\
\bottomrule
\end{tabularx}
\end{table}

% ─────────────────────────────────────────────────────────────────────────────
\section{Roman Urdu Orthographic Variation Examples}
\label{app:data-variation}

Table~\ref{tab:ortho-variation} documents the orthographic variations encountered
in the dataset for common booking-domain words.  These variations were catalogued
from real vendor conversation logs and represent the primary challenge for subword
tokenisers.

\begin{table}[H]
\centering
\caption{Roman Urdu orthographic variation examples (booking domain vocabulary).}
\label{tab:ortho-variation}
\small
\begin{tabularx}{\textwidth}{>{\raggedright\arraybackslash}p{0.18\textwidth}Y Y}
\toprule
\textbf{Urdu Word} & \textbf{Meaning} & \textbf{Observed Romanisations} \\
\midrule
waqt     & time  & waqt, vakt, wkht, waqat, tym, time \\
kal      & tomorrow / yesterday & kal, kull, cal, kl \\
raat     & night & raat, rat, raat ko, raaat, night \\
sham     & evening & sham, shaam, sham ko, evening \\
paise    & money & paise, paisa, paisay, money \\
booking  & booking & booking, buking, book karna \\
dastyab  & available & available, available hai, khali \\
\bottomrule
\end{tabularx}
\end{table}

% ─────────────────────────────────────────────────────────────────────────────
\section{Data Augmentation Statistics}
\label{app:data-augmentation}

\begin{table}[H]
\centering
\caption{Dataset size before and after augmentation, by source and class.}
\label{tab:aug-stats}
\small
\begin{tabularx}{\textwidth}{>{\raggedright\arraybackslash}p{0.34\textwidth}Yrrr}
\toprule
\textbf{Class} & \textbf{Source} & \textbf{Real} & \textbf{Augmented} & \textbf{Total} \\
\midrule
\texttt{inquiry}              & Vendor logs + manual & 120 & 157 & 277 \\
\texttt{info}                 & Manual + synthetic   &  80 &  91 & 171 \\
\texttt{transaction\_confirm} & Vendor logs + manual &  60 &  87 & 147 \\
\texttt{unknown}              & Manual + synthetic   &  20 &  42 &  62 \\
\hline
\textbf{Total}                &                      & \textbf{280} & \textbf{377} & \textbf{657} \\
\bottomrule
\end{tabularx}
\end{table}

\noindent Augmentation techniques applied (in order of application):
\begin{enumerate}
  \item \textbf{Back-translation} (English $\to$ Urdu script $\to$ Roman Urdu $\to$
        English): applied to all English utterances.
  \item \textbf{Paraphrasing via Gemini API}: 3 paraphrases per original utterance,
        filtered to remove duplicates and low-confidence outputs.
  \item \textbf{Lexical substitution}: replaced time expressions and service names
        with domain-synonyms (e.g., ``futsal'' $\leftrightarrow$ ``5-a-side'',
        ``7~PM'' $\leftrightarrow$ ``tonight'' $\leftrightarrow$ ``raat ko'').
  \item \textbf{Phonetic Roman Urdu rewriting}: applied to all Roman Urdu utterances
        using a custom rule set derived from Table~\ref{tab:ortho-variation}.
\end{enumerate}

The dataset and annotation files are available in the project repository at:\\
\url{https://github.com/JazibWaqas/JHAT/tree/main/data}.

%%% Local Variables:
%%% mode: latex
%%% TeX-master: "../report"
%%% End:
